% ECONOMICS_PROBLEM: {"id": "O14-274", "title": "Development through services", "jel_code": "O14", "source_id": "OP-0378"}
% !TeX program = xelatex
\documentclass[11pt,a4paper]{article}
\usepackage[margin=23mm]{geometry}
\usepackage{fontspec}
\setmainfont{Latin Modern Roman}
\usepackage{amsmath,amssymb,mathtools}
\usepackage{xcolor,enumitem,booktabs,longtable,array,xurl,hyperref}
\definecolor{muted}{HTML}{53636C}
\hypersetup{colorlinks=true,urlcolor=blue,linkcolor=blue}
\setlength{\parindent}{0pt}
\setlength{\parskip}{5pt}
\newcommand{\R}{\mathbb R}
\newcommand{\E}{\mathbb E}
\newcommand{\N}{\mathbb N}
\newcommand{\ind}{\mathbf 1}
\newcommand{\Var}{\operatorname{Var}}
\newcommand{\Cov}{\operatorname{Cov}}
\newcommand{\supp}{\operatorname{supp}}
\DeclareMathOperator*{\argmin}{arg\,min}
\DeclareMathOperator*{\argmax}{arg\,max}
\newcommand{\entrymeta}[3]{{\sffamily\small\color{muted}\textbf{\texttt{#1}}\quad|\quad #2\par #3\par}}
\newcommand{\Problemref}[1]{\texttt{#1}}
\newcommand{\JELref}[2]{\texttt{#2} (JEL \texttt{#1})}
\begin{document}
\section*{Development through services}
\textbf{Problem ID:} \texttt{O14-274}\quad \textbf{JEL:} \texttt{O14}\par
% BEGIN ECONOMICS STATEMENT
\label{problem:OP-0378}
{\sffamily\small\color{muted}Original subject group: Development and poverty\par}
\label{definitions:problem:30007660}
\textbf{Audit classification:} Published research agenda. Repository abstract and metadata checked. Its author-list spelling Youdon conflicts with its recommended citation; Yale Economics and NBER confirm Youdan Zhang. Baseline questions are addressed by its model; good-job criteria and extensions are editorial.
\subsection*{Definitions and precise research target}
\textbf{Complete editorial problem.} Consider a specified low-income economy with agriculture, manufacturing, and heterogeneous service industries. For sector \(s\), let \(L_{st}\) be employment, \(Y_{st}\) real value added, and, when \(L_{st}>0\), \(v_{st}=Y_{st}/L_{st}\) measured labor productivity. Let \(p\) be a feasible package affecting service technology, tradability, skills, or market access, with a declared fiscal cost. Equilibrium outcomes \(L_{st}(p)\) and \(Y_{st}(p)\) include movement of workers and demand responses. Define \(\Delta Y_T(p)=\sum_sY_{sT}(p)-\sum_sY_{sT}(0)\), and let \(\Delta J_T(p)\) measure the change in jobs meeting stated real-pay, stability, and accessibility thresholds for less-educated workers. Determine when service expansion raises both aggregate productivity and broadly accessible good employment without a preceding manufacturing expansion. Distinguish within-sector productivity growth from compositional labor reallocation and relative-price measurement effects. Identify causal responses where possible and validate structural predictions against sectoral employment and expenditure changes. Define the counterfactual technology path rather than assuming manufacturing growth is infeasible everywhere. The source poses foundational questions and proposes a baseline framework with future extensions. The quantitative feasibility criteria and employment-quality requirements here adapt that agenda rather than reproduce an unanswered theorem from the source.

\textbf{Definitions and admissible scope.} Let sectors \(s\in\{a,m,n,t\}\) denote agriculture, manufacturing, nontradable consumer services, and tradable services. Real value added uses fixed or explicitly chained prices; measured \(Y_s/L_s\) is average value added per worker-hour, not a technology parameter. A structural service-growth model specifies nonhomothetic demand, migration, sector technologies, trade costs, capital, skill endowments, and a resource/fiscal closure. Good jobs mean paid jobs meeting declared real-hourly-pay, contract-duration, and qualification thresholds, counted once per baseline worker even if several jobs meet the thresholds. The editorial question is to identify or validate interventions satisfying \(\Delta Y_T>0\) and \(\Delta J_T>0\) under the same baseline resources. A model that already computes those conditions is a partial answer, not proof of universal open status. Fix a nonnegative integer horizon \(T\), baseline adult-worker cohort, service policy menu, and resource budget. Define \(J_T(p)=\sum_i1\{i\text{ is employed at }T,\;w^{\rm real}_{iT}(p)\ge\underline w,\;\ell_{iT}(p)\ge\underline\ell,\;\text{job is accessible under declared qualifications}\}\) among baseline less-educated workers, where \(\ell_{iT}\) is the declared contract duration and pay and duration thresholds are fixed; \(\Delta J_T(p)=J_T(p)-J_T(0)\). Employment \(L\) is measured in worker-hours. \(\Delta Y_T>0\) is an output condition, not by itself a labor-productivity condition: if productivity is also required define \(V_T(p)=\sum_sY_{sT}(p)/\sum_sL_{sT}(p)\) with positive denominator and add \(V_T(p)>V_T(0)\). Service-led feasibility means no prerequisite manufacturing intervention is imposed; specify its baseline technology path, rather than treating historical sequencing as a theorem. Return the identified or model-validated feasible policy set.
\subsection*{Variants and assumption changes}
\begin{enumerate}
\item Consumer-services scope (source-motivated): restrict growth to locally nontradable retail and hospitality and analyze distributional effects under nonhomothetic demand and spatial heterogeneity. The source already supplies a baseline framework.
\item Tradable-services and access (editorial): permit service exports, skill investment, and digital-market access and impose a minimum good-job gain for less-educated workers. Compare this constrained path with manufacturing expansion under the same budget.
\end{enumerate}
\subsection*{References and source audit}
\textbf{Support and access.} Repository abstract and metadata checked. Its author-list spelling Youdon conflicts with its recommended citation; Yale Economics and NBER confirm Youdan Zhang. Baseline questions are addressed by its model; good-job criteria and extensions are editorial. \textbf{Locator:} Yale EGC DP 1124, January 2026, abstract paragraphs 2--4.
\begin{enumerate}\raggedright
\item\hypertarget{definitions:source:30007660:1}{} Michael Peters; Youdan Zhang; Fabrizio Zilibotti. 2026. \emph{Skipping the Factory: Service-Led Growth and Structural Transformation in the Developing World}. Abstract, paragraphs 2 and 4, EGC Discussion Paper 1124, January 2026.
\url{https://elischolar.library.yale.edu/egcenter-discussion-paper-series/1124/}.
DOI: \url{https://doi.org/10.3386/w34692}.
\end{enumerate}

\subsection*{Common conventions: Source mathematical conventions}

These conventions apply to each entry unless explicitly replaced.

\textbf{Models and identification.} A primitive model \(m\) specifies all probability laws, feasible choices, information, equilibrium and measurement rules used by its target. The admissible class \(\mathcal M\) is fixed independently of outcomes. If \(P_O(m)\) is its observable probability law and \(\tau(m)\) the target, its identified set at observable law \(P\) is
\[
 \mathcal I_\tau(P)=\{\tau(m):m\in\mathcal M,\ P_O(m)=P\}.
\]
Point identification means that this set is a singleton. An empty set signals incompatibility of data and assumptions. Coordinate bounds need not describe the joint set. Bounds are sharp only if every retained target is attainable or arbitrarily approachable by compatible models. Finite bounds require explicit restrictions on unobserved quantities; unrestricted confounding, hidden wealth, extrapolation or arbitrary equilibrium selection can yield uninformative sets.

\textbf{Causal interventions.} A potential outcome \(Y_i(p)\) is the outcome under a complete policy regime \(p\), including timing, financing, information and market spillovers. The target population and units are fixed before treatment. With interference, use the complete assignment vector or a justified exposure mapping; individual no-interference notation alone cannot identify an economy-wide intervention. A difference of mean potential outcomes does not require the cross-world joint distribution. Conditioning on post-treatment migration, survival, enrollment or employment changes the estimand and needs separate justification.

\textbf{Statistical statements.} An estimator, confidence set, forecast or test specifies a sampling law, model class and loss/coverage criterion. Uniform \((1-\alpha)\)-coverage means \(\inf_{m\in\mathcal M}\Pr_m\{\tau(m)\in C_n\}\geq1-\alpha\), or an explicitly asymptotic version. The whole parameter space trivially has coverage, so informativeness requires a stated width, risk or power objective. A forecasting improvement is relative to a named benchmark, information set and evaluation population.

\textbf{Equilibrium and optimization.} An equilibrium comprises feasible plans, optimality, beliefs where needed, budget constraints and all market/resource clearing. The primitives or a stated benchmark define each correspondence \(\mathcal E(p;m)\). Nonemptiness is assumed or proved, never inferred just from compact policy sets. A selected-equilibrium target uses a declared selection rule; a robust target quantifies over all permitted equilibria. Use suprema or infima unless attainment is established. An \(\varepsilon\)-optimal policy is feasible and within \(\varepsilon\) of the finite optimum value. Report an empty feasible set separately.

\textbf{Welfare and accounting.} Normative utility, interpersonal normalization, social weights, numeraire, discounting and the use of public receipts are primitives. Behavioral utility need not coincide with the welfare criterion. Household welfare includes ownership income and rents once; adding profits again to compensating variation generally double-counts them. A monetary equivalent is defined by an explicit transfer and price convention, with monotonicity and range conditions. Average effects, mechanism contributions, transfers and welfare losses are different targets. Infinite sums require convergence and dynamic feasible plans require terminal or no-Ponzi conditions.

\textbf{Source versions.} A working-paper date, revision date and journal publication year are distinguished. A DOI for a journal version is not automatically the DOI of an earlier manuscript. Locators refer to the stated version and printed pages where specified. A metadata-only check does not verify a claim about a full-text conclusion. Every entry's source subsection narrows the provenance claim accordingly.
% END ECONOMICS STATEMENT
\end{document}
